\begin{lemma}[Finite dependence on the enumeration orbit]
For fixed $r,n,k$, the copied value
$\mathsf{PowerCopiedLeft}_r(h,x,n,k)$ depends only on the finitely many
values of $h$ on the first $n+k+2$ shifts of $x$.  Thus, if $x$ and $y$
give the same $h$-values at all these shifts, then
\[
\mathsf{PowerCopiedLeft}_r(h,x,n,k)
=\mathsf{PowerCopiedLeft}_r(h,y,n,k).
\]
\end{lemma}
\begin{proof}
We prove the following stronger assertion by induction on the depth $m$:
for every row $i$, if
\[
 h(\mathsf{PowerShiftIter}(j,x))
 =h(\mathsf{PowerShiftIter}(j,y))
 \qquad\text{for every }j\leq i+m+2,
 \tag{1}
\]
then
\[
 \mathsf{PowerCopiedLeft}_r(h,x,i,m)
 =\mathsf{PowerCopiedLeft}_r(h,y,i,m).
 \tag{2}
\]

For the base case $m=0$, fix an arbitrary row $i$ and assume (1), so in
particular the assumed agreement holds at $j=i$ and at $j=i+1$, since
$i\leq i+2$ and $i+1\leq i+2$.  The depth-zero clause of
\noderef{PowerCopiedLeft}, instantiated at $i$, gives
\[
\begin{aligned}
\mathsf{PowerCopiedLeft}_r(h,x,i,0)
  &=\mathsf{PowerIResponse}_r\bigl(
      h(\mathsf{PowerShiftIter}(i,x)),
      h(\mathsf{PowerShiftIter}(i+1,x))\bigr),\\
\mathsf{PowerCopiedLeft}_r(h,y,i,0)
  &=\mathsf{PowerIResponse}_r\bigl(
      h(\mathsf{PowerShiftIter}(i,y)),
      h(\mathsf{PowerShiftIter}(i+1,y))\bigr).
\end{aligned}
\tag{3}
\]
The two arguments in the first response in (3) equal the corresponding
arguments in the second response by the cases $j=i$ and $j=i+1$ of (1).
By the definition of the fixed response function
\noderef{PowerIResponse}, equal pairs of inputs have equal responses.
Thus (2) holds for this arbitrary $i$, and hence for every row at depth
zero.

For the successor step, assume the strengthened assertion at depth $m$.
Fix an arbitrary row $i$ and assume agreement through the successor bound
\[
 j\leq i+(m+1)+2,
 \qquad\text{equivalently }j\leq i+m+3.
 \tag{4}
\]
First, if $j\leq i+m+2$, then $j\leq i+m+3$ by monotonicity of addition,
so (4) supplies the induction hypothesis at depth $m$ for row $i$:
\[
 \mathsf{PowerCopiedLeft}_r(h,x,i,m)
 =\mathsf{PowerCopiedLeft}_r(h,y,i,m).
 \tag{5}
\]
Second, the same full bound (4) is the bound needed for row $i+1$, because
\[
 (i+1)+m+2=i+m+3.
\]
Consequently the induction hypothesis at depth $m$ for row $i+1$ gives
\[
 \mathsf{PowerCopiedLeft}_r(h,x,i+1,m)
 =\mathsf{PowerCopiedLeft}_r(h,y,i+1,m).
 \tag{6}
\]
The successor-depth clause of \noderef{PowerCopiedLeft}, applied at row
$i$ for the two enumerations, gives
\[
\begin{aligned}
\mathsf{PowerCopiedLeft}_r(h,x,i,m+1)
 &=\mathsf{PowerIResponse}_r\bigl(
     \mathsf{PowerCopiedLeft}_r(h,x,i,m),
     \mathsf{PowerCopiedLeft}_r(h,x,i+1,m)\bigr),\\
\mathsf{PowerCopiedLeft}_r(h,y,i,m+1)
 &=\mathsf{PowerIResponse}_r\bigl(
     \mathsf{PowerCopiedLeft}_r(h,y,i,m),
     \mathsf{PowerCopiedLeft}_r(h,y,i+1,m)\bigr).
\end{aligned}
\tag{7}
\]
The first inputs in (7) are equal by (5), and the second inputs are equal
by (6).  Therefore the two pairs supplied to the same deterministic
function \noderef{PowerIResponse} are equal, so the two outputs in (7) are
equal.  This proves the strengthened assertion at depth $m+1$ for the
arbitrary row $i$, and hence for every row.

Finally instantiate the strengthened assertion at depth $k$ and row
$i=n$.  Its premise then says that for every $j\leq n+k+2$,
\[
h(\mathsf{PowerShiftIter}(j,x))
=h(\mathsf{PowerShiftIter}(j,y)),
\]
which is exactly the stated hypothesis \(\mathsf{hag}\); its conclusion is
the required equality.
\end{proof}
